\section{Tổng quan framework DeSFAM}
\label{sec:desfam}

\subsection{Kiến trúc ba mô-đun}
DeSFAM (\emph{Dynamic eBPF-Driven Syscall Filtering and Anomaly
Mitigation})~\cite{desfam_original} là một framework bảo mật thời gian thực, nhẹ cho
container, thực thi nguyên tắc \emph{đặc quyền tối thiểu} ở tầng syscall nhằm giảm
thiểu leo thang đặc quyền, thoát container và chèn syscall. DeSFAM hoạt động theo ba
mô-đun (giai đoạn) tích hợp (Hình~\ref{fig:arch}): lập hồ sơ/danh sách truy cập
syscall ở không gian người dùng, phát hiện bất thường, và thực thi trong nhân.

\begin{figure}[htbp]
  \centering
  \includegraphics[width=\linewidth]{figures/desfam-architecture}
  \caption{Kiến trúc ba mô-đun của DeSFAM. Báo cáo này đặt trọng tâm vào Mô-đun~2
           (SyscallAD)---phát hiện bất thường syscall; Mô-đun~1 và Mô-đun~3 chỉ
           được điểm qua.}
  \label{fig:arch}
\end{figure}

\subsection{Mô-đun 1 --- Lập danh sách truy cập syscall lai và hồ sơ seccomp (điểm qua)}
\label{subsec:m1}
Mô-đun~1 sinh một danh sách truy cập syscall tối thiểu, dành riêng cho từng
container, bằng cách kết hợp ba nguồn: (i) \emph{phân tích tĩnh}---dịch ngược binary
ELF (\texttt{objdump}, \texttt{readelf}) và dựng đồ thị lời gọi để rút tập syscall
tĩnh $S_\text{static}$; (ii) \emph{lập hồ sơ động}---gắn eBPF probe dưới các khối
lượng công việc chuẩn để thu tập syscall động $S_\text{dynamic}$; (iii) \emph{lọc
theo CVE}---trích tên syscall rủi ro cao từ mô tả CVE (NVD) bằng NLP để tạo
$S_\text{blocked}$. Danh sách cuối cùng được biên dịch thành hồ sơ seccomp và nạp
động qua \texttt{policy\_loader}. \emph{Mô-đun này nằm ngoài trọng tâm báo cáo và
không được hiện thực hóa lại}; ở đây chỉ ghi nhận rằng nó cung cấp lớp phòng thủ
tĩnh, bổ trợ cho lớp phát hiện động của Mô-đun~2.

\subsection{Mô-đun 3 --- Thực thi thích ứng và quản lý chính sách (điểm qua)}
\label{subsec:m3}
Mô-đun~3 nhận các chuỗi bị gắn cờ từ SyscallAD và quyết định hành động bằng
\emph{điểm rủi ro} tổng hợp từ: ánh xạ syscall sang kỹ thuật MITRE ATT\&CK, mức
nghiêm trọng tích lũy, và khớp mẫu khai thác CVE. Theo ngưỡng rủi ro, hệ thống
\emph{chặn} / \emph{giới hạn tốc độ} / \emph{ghi nhật ký}; việc chặn được thực thi
trong nhân bằng eBPF gắn vào LSM hook (\texttt{bpf\_override\_return(-EPERM)}). DeSFAM
còn có vòng tự phục hồi cập nhật lại danh sách seccomp khi một syscall bị chặn lặp
lại. \emph{Mô-đun này cũng nằm ngoài trọng tâm}: báo cáo dừng ở lớp phát hiện, còn
việc chấm điểm rủi ro và chặn nội tuyến được nêu là hướng mở rộng.

\subsection{Vị trí của SyscallAD và lý do chọn làm trọng tâm}
SyscallAD (Mô-đun~2) là \emph{trái tim phát hiện} của DeSFAM: Mô-đun~1 chỉ thu hẹp
bề mặt tấn công theo chính sách tĩnh, còn Mô-đun~3 chỉ \emph{phản ứng} dựa trên cờ
mà SyscallAD phát ra. Nói cách khác, năng lực ``nhận biết tấn công'' của cả
framework hội tụ ở SyscallAD. Đây cũng là phần mang tính \emph{học máy} và \emph{khả
chuyển} nhất---độc lập với cơ chế thực thi của một nhân cụ thể---nên phù hợp nhất để
nghiên cứu, tái lập và đánh giá định lượng. Vì vậy báo cáo dành toàn bộ
Mục~\ref{sec:syscallad} và~\ref{sec:privesc} cho mô-đun này.
